% !TeX root = ../main.tex
%------------------------------------------------------------------------------
%  Engineering Geology
%  Sources: Hoek & Brown (1980) Ch. 3; Brady & Brown (2006) Ch. 3;
%  Knill (2003) on the ground model.
%------------------------------------------------------------------------------
\chapter{Engineering Geology}
\label{ch:engineering-geology}

The parameters used in underground design are abstractions of geological
conditions. Intact strength stands for mineralogy, fabric and cementation; the
Geological Strength Index stands for the arrangement of discontinuities; the
disturbance factor stands for the damage excavation inflicts on both. These
abstractions are necessary, but each discards information, and a parameter
applied without knowing what was discarded is liable to be used outside its
range. This chapter establishes the geological basis on which the
parameters of Chapters~\ref{ch:rock-mechanics}
and~\ref{ch:rock-mass-parameters} rest, and argues that the feature governing
an underground project is more often geological than mechanical: a weathering
profile, a fault zone or a bedding plane at an unfavourable orientation. Later
analytical refinement does not compensate for a failure to identify such a
feature.

\section{Ground materials and structure}
\label{sec:ground-materials}

\subsection{The geological model and the ground model}
\label{sec:ground-model}

Three distinct objects are routinely called the same thing. The
\emph{geological model} is an interpretation of what the ground is and how it
came to be: units, contacts, structures and their history. The \emph{ground
model} attaches engineering meaning to that interpretation: material
behaviour, parameters, groundwater, and the uncertainty in each. The
\emph{design model} is the idealisation actually analysed, simplified to what
the calculation can accommodate. Chapter~\ref{ch:ground-modelling} keeps
these three names and adds a fourth, the \emph{hydrogeological model}, which
carries aquifers, heads and flow boundaries (Section~\ref{sec:model-types}).

Each is a different kind of statement.
Disagreement about the geological model is resolved by more or better
observation. Disagreement about the ground model is resolved by testing and
correlation. Disagreement about the design model is resolved by argument about
what may safely be neglected. Conflating them produces the familiar situation
in which a modelling assumption is defended as though it were a geological
fact.

The ground model, rather than the geological interpretation, is the central
engineering deliverable of an investigation. It is the object handed to the
designer, updated during construction and passed to the operator, and it
should be built, versioned and owned accordingly
(Chapter~\ref{ch:ground-modelling}).

\subsection{Rock materials and weathering}
\label{sec:rock-materials}

Rocks are usefully grouped by what they do in an excavation rather than by
their origin alone. Igneous intrusives are typically strong, isotropic and
jointed in regular sets, so at stress levels low relative to their intact
strength their behaviour is structurally controlled, and their difficulties
concentrate at contacts and in weathered zones; at high stress the same mass
becomes stress-controlled (Chapter~\ref{ch:high-stress}).
Sedimentary sequences are layered, and bedding introduces both anisotropy and
a set of persistent, often clay-coated, discontinuities that frequently govern
roof stability. Metamorphic rocks with strong foliation (slates, schists,
phyllites) are the most awkward, combining marked strength anisotropy with
a pervasive weakness plane whose orientation relative to the excavation
decides the outcome.

Weathering modifies all of them, and its variation with depth matters most
for shallow works. A weathering profile grades from residual soil
through completely, highly, moderately and slightly weathered rock to fresh
rock, and
the transition is rarely uniform: it follows joints and faults downward,
leaving corestones of fresh rock in a weathered matrix and producing a
rockhead surface far more irregular than a borehole spacing of tens of metres
can resolve (Figure~\ref{fig:geo-weathering-profile}).

Two errors follow from treating weathering grade as a proxy for strength. The
first is assuming the correlation holds across a lithological change; where a
formation contains both a weak and a strong member, a more weathered specimen
of the strong member can return a higher intact strength than a fresher
specimen of the weak one. The
second is assuming the profile is layered. An irregular rockhead means that a
tunnel nominally in rock may meet a soil-filled trough over part of its
length, and mixed-face conditions are among the most damaging surprises in
shallow tunnelling (Figure~\ref{fig:r3d-weathering-block}).

\begin{figure}[htbp]
\centering
\datafigure[\textwidth]{ch04/geo-weathering-profile}
\caption{Schematic weathering profile, grades VI (residual soil) to I (fresh)
marked at right. Weathering descends along a fault (dashed), leaving
corestones in grades IV and III. Boreholes 40~m apart all find rockhead (heavy
line) above the 6~m tunnel, which meets soil-like ground over 18~m of its
length (hatched). Vertical scale exaggerated.}
\label{fig:geo-weathering-profile}
\end{figure}

\begin{figure}[htbp]
\centering
\renderfigure{weathering_block}{%
  \rlab[south]{0.33,0.49}{0.08,0.86}{Mixed face}
  \rlab[south]{0.465,0.63}{0.42,0.86}{Borehole}
  \rlab[south]{0.70,0.47}{0.70,0.84}{Corestones}
  \rlab[north]{0.375,0.30}{0.31,0.12}{Fault}
  \rlab[north]{0.62,0.22}{0.62,0.06}{Fresh rock}
}
\caption{The profile of Figure~\ref{fig:geo-weathering-profile} at true scale,
extended 30~m along strike (illustrative) and cut on the tunnel axis: grades VI
to I from light to dark grey, the fault (carmine tint) carrying weathering down
into a trough. The three boreholes miss the trough; the lining is carmine
over the 18~m where it meets soil-like ground.}
\label{fig:r3d-weathering-block}
\end{figure}

\subsection{Discontinuities and structure}
\label{sec:discontinuities}

At excavation scale, the behaviour of a rock mass is governed largely by its
discontinuities. Their
description is standardised (orientation, spacing, persistence, roughness,
aperture, infilling, seepage and the number of sets), and each of these
parameters feeds a specific downstream method. Orientation and set count
determine block formation (Chapter~\ref{ch:structural-analysis}); spacing and
persistence determine block size; roughness, aperture and infilling determine
shear strength; seepage and infilling condition enter the classification
systems (Chapter~\ref{ch:classification}).

Discontinuity shear strength links this description to mechanics. A
rough unfilled joint mobilises strength above its basic friction angle through
the geometric effect of asperities, and that additional component reduces as
normal stress rises and asperities are sheared through. This is the bilinear description of
\citet{patton1966multiple}, refined by \citet{barton1977shear} into a continuous relation in terms
of joint roughness and wall strength (the relation is given with its
constants in Appendix~\ref{app:constants}, and drawn in
Figure~\ref{fig:geo-joint-shear}). Two consequences carry into design.
Roughness is scale-dependent, so a value measured on a 100\,mm laboratory
specimen overestimates the strength of a joint several metres across. An
infilling thicker than the asperity amplitude decouples the surfaces
altogether: the strength becomes that of the infill, which may be a fraction
of the rock-on-rock value (Section~\ref{sec:discontinuities-and-soils}).

\begin{figure}[htbp]
\centering
\datafigure[\textwidth]{ch04/geo-joint-shear}
\caption{Left: shear strength of a rough unfilled joint, $\phi_b = 30^\circ$,
JCS~=~60~MPa, from the Barton relation for three values of JRC, with the
bilinear envelope for a $15^\circ$ asperity angle sheared through at
1.5~MPa. Right: the roughness term for JRC~10 on a fresh wall and a weathered
one at a quarter of the strength; it falls with normal stress and vanishes at JCS.}
\label{fig:geo-joint-shear}
\end{figure}

Structure at alignment scale is dominated by faults. A fault is a zone
rather than a plane, with a core of gouge or breccia and a surrounding damage zone of
elevated fracture density (Figure~\ref{fig:geo-fault-zone}), and both are usually weaker and more permeable than
the host rock. Fault crossings therefore combine reduced strength, increased
deformability, elevated inflow and abrupt change in ground class within a few
metres of advance. They occupy a small fraction of most alignments and a large
fraction of most risk registers, and locating them is among the highest-value
objectives of an investigation.

\begin{figure}[htbp]
\centering
\datafigure[\textwidth]{ch04/geo-fault-zone}
\caption{Schematic of a fault crossed by a tunnel (plan, tunnel between the
two lines), with fracture density, permeability and strength relative to the
host rock. A 6~m core lies within a damage zone some 50~m across; properties
change through the damage zone and step at the core. Values illustrative.}
\label{fig:geo-fault-zone}
\end{figure}

\section{Stress history and problem ground}
\label{sec:stress-history-problem-ground}

The present state of a rock mass records its history, and two aspects of that
history bear directly on design. Erosional unloading leaves horizontal
stresses that are high relative to the vertical stress, because horizontal
strain was restrained while overburden was removed; this is the principal
reason shallow rock excavations can meet horizontal-to-vertical stress ratios
well above unity, treated quantitatively in
Chapter~\ref{ch:in-situ-stress}. Where the stress is high relative to the
intact strength, core recovered from a borehole breaks into discs, and core
discing is the first field indication of it. Overconsolidated clays similarly retain
lateral stresses from a greater former overburden, with consequences for
excavation support and for heave.

Locked-in residual stress at grain and block scale is a smaller effect but
explains observations that otherwise appear anomalous: time-dependent
deformation of specimens after recovery, and stress measurements that scatter
beyond their instrumental precision.

Certain ground types are difficult enough to deserve identification during the
geological study rather than discovery during construction. Karst carbonates
carry solution cavities and conduits that can be metres across, may be
air-filled, water-filled or sediment-filled, and are essentially invisible
between boreholes. Swelling ground (smectite-rich mudrocks, and anhydrite
converting to gypsum with a volume increase of roughly 60 per cent on
hydration) generates invert heave
that continues for years and can lift a completed lining. Gas-bearing strata,
particularly coal measures and organic-rich sediments, introduce methane and
an explosion risk that changes the construction regime. Ultra-weak and
heterogeneous rock masses such as tectonised flysch defeat both intact rock
testing and conventional classification, because the parameter that matters is
the proportion and connectivity of the weak member rather than the strength of
either component.

Each of these has a geological signature that a competent desk study will
find: a carbonate formation with known karst development in the region, a
mudrock with a recorded swelling history, a coal seam within the influence
zone. Identifying them at desk-study stage costs little compared with meeting
them at the face. Table~\ref{tab:problem-ground} sets each ground type against
its signature and its consequence.

\begin{table}[htbp]
\caption{Problem ground: the geological signature a desk study can find and the consequence for the excavation.}
\label{tab:problem-ground}
\centering
\small
\begin{tabularx}{\textwidth}{>{\raggedright\arraybackslash}p{0.21\textwidth} Y Y}
\toprule
\textbf{Ground type} & \textbf{Geological signature} & \textbf{Consequence} \\
\midrule
Karst carbonates & Carbonate formation with known karst development in the region & Solution cavities and conduits up to metres across, filled with air, water or sediment, and essentially invisible between boreholes \\
Swelling ground & Mudrock with a recorded swelling history; smectite-rich mudrocks; anhydrite & Invert heave that continues for years and can lift a completed lining \\
Gas-bearing strata & Coal seam within the influence zone; coal measures and organic-rich sediments & Methane and an explosion risk that changes the construction regime \\
Ultra-weak and heterogeneous rock masses & Tectonised sequence such as flysch & Intact rock testing and conventional classification do not apply; the proportion and connectivity of the weak member govern \\
\bottomrule
\end{tabularx}
\end{table}

\section{Hazards, uncertainty and baselining}
\label{sec:geological-hazards}

The engineering output of the geological study is a hazard map along the
alignment: where the ground changes, where structure is expected, where water
is expected, and where the consequence of being wrong is highest. That map
drives investigation targeting (Chapter~\ref{ch:site-investigation}), and it should also drive alignment and depth
selection, which are the cheapest available means of avoiding a hazard.

Geological uncertainty is of three kinds. Some of it is
\emph{positional}: the unit is correctly identified but its boundary is not
located precisely. Some is \emph{interpretive}: the observations admit more
than one geological explanation, and the alternatives have different
engineering consequences. Some is \emph{parametric}: the unit and its geometry
are agreed, but its properties are variable or poorly sampled. Investigation
reduces the first two efficiently and the third only slowly, because one
well-placed observation can locate a boundary or discriminate between two
interpretations outright, whereas sampling scatter falls only as the square
root of the number of specimens and never below the variability of the
material itself. The distinction determines what additional drilling can be
expected to achieve (Figure~\ref{fig:geo-uncertainty}).

\begin{figure}[htbp]
\centering
\datafigure[0.62\textwidth]{ch04/geo-uncertainty}
\caption{Schematic of the three kinds of uncertainty against well-placed
observations: positional halves with each hole that brackets a boundary,
interpretive collapses at a discriminating observation (here the third), and
parametric at an unsampled point, $\propto\sqrt{1 + 1/n}$, falls only towards
the material's own scatter (dotted).}
\label{fig:geo-uncertainty}
\end{figure}

Interpretive uncertainty is the type most often suppressed. A geological section is drawn as one interpretation,
and once drawn it acquires an authority the data do not support. Where the
observations admit two structural interpretations (a fault or a
facies change, a continuous or a discontinuous marker), both should be
carried forward, with their differing engineering consequences stated, until
evidence distinguishes them.

Finally, the ground model becomes a contractual object. A geotechnical
baseline report states the ground conditions the contract assumes, in terms
that can be tested against what is encountered, and it necessarily
simplifies the ground model into baselined statements: this many
metres of each ground class, this range of water inflow, faults of this
character. Writing those statements is an engineering geological task with
commercial consequences, and it is treated in
Chapter~\ref{ch:delivery}. A baseline can be only as testable as the
ground model behind it is explicit.

\chapterreferences
